\subsection{Đề 2}
\Opensolutionfile{ans}[ans/ans-2-B9-De2-LC]
\caulc
\begin{ex}%[2D3N1-1]
	Khoảng tứ phân vị $\Delta Q$ của một mẫu số liệu ghép nhóm là
	\choice
	{\True Hiệu giữa tứ phân vị thứ ba và tứ phân vị thứ nhất của mẫu số liệu đó: $\Delta Q=Q_3-Q_1$}
	{Hiệu giữa tứ phân vị thứ nhất và tứ phân vị thứ ba của mẫu số liệu đó: $\Delta Q=Q_1-Q_3$}
	{Tổng giữa tứ phân vị thứ nhất và tứ phân vị thứ ba của mẫu số liệu đó: $\Delta Q=Q_1+Q_3$}
	{Trung bình cộng của tứ phân vị thứ nhất và tứ phân vị thứ ba của mẫu số liệu đó: $\Delta Q=\dfrac{Q_1+Q_3}{2}$}
	\loigiai{
	Theo định nghĩa ta có: Khoảng tứ phân vị $\Delta Q$ của một mẫu số liệu ghép nhóm là hiệu giữa tứ phân vị thứ ba và tứ phân vị thứ nhất của mẫu số liệu đó: $\Delta Q=Q_3-Q_1$.
}
\end{ex}
\begin{ex}%[2D3H1-3]
	Bác Long làm nghề xe ôm, thống kê quãng đường đi được của bác Long trong $20$ ngày được cho trong bảng sau
	\begin{center}
		\begin{tabular}{|c|c|c|c|c|c|}
		\hline
	Quãng đường (km)	&$\left[ 20;30\right)$  &$\left[ 30;40\right)$  & $\left[ 40;50\right)$ & $\left[ 50;60\right)$ & $\left[ 60;70\right)$ \\
		\hline
		Số ngày&$3$  &$5  $&$6$  &$4$  & $2$ \\
		\hline
	\end{tabular}
	\end{center}
Khoảng tứ phân vị của mẫu số liệu trên là
	\choice
	{\True $18{,}5$}
	{$20{,}5$}
	{$26{,}5$}
	{$22{,}5$}
	\loigiai{
		Cỡ mẫu: $n=20$.\\
		Gọi $x_1; x_2;\ldots ; x_{20}$ là mẫu số liệu gốc gồm quãng đường đi được của bác Long được sắp xếp thứ tự không giảm.\\
		Ta có $x_1; x_2; x_3 \in \left[ 20;30\right)$ \\
		$x_4; x_5; x_6; x_7; x_8 \in \left[ 30;40\right)$\\
		$x_9; x_{10}; x_{11}; x_{12}; x_{13}; x_{14} \in \left[ 40;50\right)$\\
		$x_{15}; x_{16}; x_{17}; x_{18} \in \left[ 50;60\right)$\\
		$x_{19}; x_{20} \in \left[ 60;70\right)$\\
		Tứ phân vị thứ nhất của mẫu số liệu gốc là $\dfrac{1}{2}\left(x_5+x_6\right) \in \left[ 30;40\right)$. Do đó, tứ phân vị thứ nhất của mẫu số liệu ghép nhóm là
		\begin{align*}
			Q_1= 30+\dfrac{\dfrac{20}{4}-3}{5}\cdot 10=34.
		\end{align*}
		Tứ phân vị thứ ba của mẫu số liệu gốc là $\dfrac{1}{2}\left(x_{15}+x_{16}\right) \in \left[50;60\right)$. Do đó, tứ phân vị thứ ba của mẫu số liệu ghép nhóm là
		\begin{align*}
			Q_3= 50+\dfrac{3\cdot \dfrac{20}{4}-14}{4}\cdot 10= 52{,}5.
		\end{align*}
	Vậy khoảng tứ phân vị cần tìm là: $\Delta Q=Q_3-Q_1=52{,}5-34=18{,}5$.
	}
\end{ex}

\begin{ex}%[2D3N1-2]
	Thống kê số học sinh đi học muộn trong một tháng của trường THPT A được cho trong bảng sau
	\begin{center}
		\begin{tabular}{|c|c|c|c|c|c|}
			\hline
			số học sinh đi muộn	&$\left[ 0;3\right)$  &$\left[ 3;6\right)$  & $\left[ 6;9\right)$ & $\left[ 9;12\right)$ & $\left[ 12;15\right)$ \\
			\hline
			Số ngày&$3$  &$5  $&$6$  &$4$  & $2$ \\
			\hline
		\end{tabular}
	\end{center}
	Khoảng biến thiên của mẫu số liệu trên là
	\choice
	{ $8$}
	{$14$}
	{$12$}
	{\True $15$}
	\loigiai{
		Khoảng biến thiên của mẫu số liệu trên là $15-0=15$.
	}
\end{ex}
\begin{ex}%[2D3H1-3]
	Anh Minh làm nghề giao hàng, thống kê quãng đường đi được của anh Minh trong $15$ ngày được cho trong bảng sau:
	\begin{center}
		\begin{tabular}{|c|c|c|c|c|c|}
			\hline
			Quãng đường (km) & $\left[ 10;20\right)$ & $\left[ 20;30\right)$ & $\left[ 30;40\right)$ & $\left[ 40;50\right)$ & $\left[ 50;60\right)$ \\
			\hline
			Số ngày & $3$ & $4$ & $5$ & $2$ & $1$ \\
			\hline
		\end{tabular}
	\end{center}
	Khoảng tứ phân vị của mẫu số liệu trên là
	\choice
	{$8$}
	{$10{,}5$}
	{\True $12{,}5$}
	{$14$}
	\loigiai{
		Cỡ mẫu: $n=15$.\\
		Gọi $x_1, x_2, \ldots, x_{15}$ là mẫu số liệu gốc gồm quãng đường đi được của anh Minh được sắp xếp theo thứ tự không giảm.\\
		Ta có $x_1, x_2, x_3 \in \left[ 10;20\right)$ \\
		$x_4, x_5, x_6, x_7 \in \left[ 20;30\right)$ \\
		$x_8, x_9, x_{10}, x_{11}, x_{12} \in \left[ 30;40\right)$ \\
		$x_{13}, x_{14} \in \left[ 40;50\right)$ \\
		$x_{15} \in \left[ 50;60\right)$ \\
		Tứ phân vị thứ nhất của mẫu số liệu gốc là $\dfrac{1}{2}(x_4 + x_5) \in \left[ 20;30\right)$. Do đó, tứ phân vị thứ nhất của mẫu số liệu ghép nhóm là
		\begin{align*}
			Q_1 = 20 + \left(\dfrac{\dfrac{15}{4} - 3}{4}\right) \cdot 10 = 21{,}25.
		\end{align*}
		Tứ phân vị thứ ba của mẫu số liệu gốc là $\dfrac{1}{2}(x_{12} + x_{13}) \in \left[ 40;50\right)$. Do đó, tứ phân vị thứ ba của mẫu số liệu ghép nhóm là
		\begin{align*}
			Q_3 = 40 + \left(\dfrac{3 \cdot \dfrac{15}{4} - 12}{2}\right) \cdot 10 = 33{,}75.
		\end{align*}
		Vậy khoảng tứ phân vị cần tìm là: $\Delta Q = Q_3 - Q_1 = 33{,}75 - 21{,}25 = 12{,}5$.
	}
\end{ex}

\begin{ex}
	Phát biểu nào sau đây là \textbf{sai}?
	\choice
	{Khoảng tứ phân vị càng nhỏ thì mẫu số liệu càng tập trung}
	{Khoảng tứ phân vị ít bị ảnh hưởng bởi các giá trị bất thường}
	{Khoảng biến thiên càng lớn thì độ phân tán càng lớn}
	{\True Khoảng tứ phân vị phụ thuộc vào các giá trị bất thường}
	\loigiai{
		Khoảng tứ phân vị $\Delta Q$ là thước đo độ phân tán của tập dữ liệu và ít bị ảnh hưởng bởi các giá trị bất thường. Nó được tính bằng hiệu giữa tứ phân vị thứ ba ($Q_3$) và tứ phân vị thứ nhất ($Q_1$).
	}
\end{ex}
%%%%%%
\begin{ex}%[2D3H1-3]
	Bảng sau thống kê lại cân nặng của $50$ quả dưa hấu được lựa chọn ngẫu nhiên sau khi thu hoạch ở nông trường:
	\begin{center}
		\begin{tabular}{|c|c|c|c|c|c|}
			\hline
			Cân nặng (kg) & $\left[ 1{,}5;1{,}6\right)$ & $\left[ 1{,}6;1{,}7\right)$ & $\left[ 1{,}7;1{,}8\right)$ & $\left[ 1{,}8;1{,}9\right)$ & $\left[ 1{,}9;2{,}0\right)$ \\
			\hline
			Số quả dưa & $5$ & $10$ & $20$ & $10$ & $5$ \\
			\hline
		\end{tabular}
	\end{center}
	Khoảng tứ phân vị của mẫu số liệu trên là
	\choice
	{\True $0{,}15$}
	{$0{,}30$}
	{$0{,}25$}
	{$0{,}40$}
	\loigiai{
		Cỡ mẫu: $n=50$.\\
		Gọi $x_1, x_2, \ldots, x_{50}$ là mẫu số liệu gốc gồm $50$ quả dưa hấu được sắp xếp thứ tự không giảm.\\
		Ta có $x_1, x_2, \ldots, x_5 \in \left[ 1{,}5;1{,}6\right)$ \\
		$x_6, x_7, \ldots, x_{15} \in \left[ 1{,}6;1{,}7\right)$ \\
		$x_{16}, x_{17}, \ldots, x_{35} \in \left[ 1{,}7;1{,}8\right)$ \\
		$x_{36}, x_{37}, \ldots, x_{45} \in \left[ 1{,}8;1{,}9\right)$ \\
		$x_{46}, x_{47}, \ldots, x_{50} \in \left[ 1{,}9;2{,}0\right)$ \\
		Tứ phân vị thứ nhất của mẫu số liệu gốc là $x_{13} \in \left[1{,}6;1{,}7\right)$. Do đó, tứ phân vị thứ nhất của mẫu số liệu ghép nhóm là
		\begin{align*}
			Q_1 = 1{,}6 + \dfrac{\dfrac{50}{4} - 5}{10} \cdot (1{,}7 - 1{,}6) = 1{,}6 + \dfrac{12{,}5 - 5}{10} \cdot 0{,}1 = 1{,}6 + \dfrac{7{,}5}{10} \cdot 0{,}1 = 1{,}6 + 0{,}075 = 1{,}675.
		\end{align*}
		Tứ phân vị thứ ba của mẫu số liệu gốc là $x_{38} \in \left[1{,}8;1{,}9\right)$. Do đó, tứ phân vị thứ ba của mẫu số liệu ghép nhóm là
		\begin{align*}
			Q_3 = 1{,}8 + \dfrac{3 \cdot \dfrac{50}{4} - (5 + 10 + 20)}{10} \cdot (1{,}9 - 1{,}8) = 1{,}8 + \dfrac{2{,}5}{10} \cdot 0{,}1 = 1{,}8 + 0{,}025 = 1{,}825.
		\end{align*}
		Vậy khoảng tứ phân vị cần tìm là: $\Delta Q = Q_3 - Q_1 = 1{,}825 - 1{,}675 = 0{,}15$.
	}
\end{ex}
\begin{ex}%[2D3V1-3]
	\immini{
	Cho bảng biểu diễn mẫu số liệu ghép nhóm về chiều cao của $50$ mẫu cây ở một vườn thực vật (đơn vị centimet). Mệnh đề nào sau đây đúng.}
{
		\begin{tabular}{|c|c|c|}
			\hline Nhóm & Tần số & Tần số tích luỹ \\
			\hline $[40; 45)$ & $7$ & $7$ \\
			$[45; 50)$& $12$ & $19$ \\
			$[50; 55)$& $9$ & $28$ \\
			$[55; 60)$& $10$ & $38$ \\
			$[60; 65)$& $8$ & $46$ \\
			$[65; 70)$& $4$ & $50$ \\
			\hline & $n=50$ & \\
			\hline
		\end{tabular}}
	\choice
	{Cỡ mẫu $n=51$}
	{$Q_1=45$}
	{\True $Q_3=12{,}65$}
	{Nhóm có tần số lớn nhất là nhóm $[55; 60)$}
	\loigiai{
		Số phần tử của mẫu là $n=50$.\\
		Ta có $\dfrac{n}{4}=\dfrac{42}{4}=10{,}5$ mà $5<10{,}5<15$.
		$$Q_1=45+\left(\dfrac{10{,}5-5}{10}\right) \cdot 5=47{,}75~(\mathrm{cm}).$$
		Ta có $\dfrac{3 n}{4}=\dfrac{3 \cdot 42}{4}=31{,}5$ mà $31<31{,}5<38$.\\
		Áp dụng công thức, ta có tứ phân vị thứ ba là
		$$Q_3=60+\left(\dfrac{31{,}5-31}{7}\right) \cdot 5 \approx 60{,}4~(\mathrm{cm}).$$
		Vậy khoảng tứ phân vị của mẫu số liệu ghép nhóm đã cho là
		\[
		\Delta_Q=Q_3-Q_1 \approx 60{,}4-47{,}75 \approx 12{,}65~(\mathrm{cm}).
		\]
	}
\end{ex}

\begin{ex}%[2D3N1-4]
	Trong một cuộc khảo sát về thói quen đọc sách của $50$ học sinh, số lượng sách mà mỗi học sinh đọc trong một tháng được thống kê như sau:
	\begin{center}
		\begin{tabular}{|c|c|c|}
			\hline
			Số lượng sách & Số học sinh & Tần số tích luỹ \\
			\hline
			$[0; 1)$ & $8$ &$8$\\
			$[1; 2)$ & $12$ &$20$\\
			$[2; 3)$ & $15$ &$35$\\
			$[3; 4)$ & $10$& $45$\\
			$[4; 5)$ & $5$ &$50$\\
			\hline
		\end{tabular}
	\end{center}
	Phần lớn số học sinh đọc được số lượng cuốn sách trong khoảng
	\choice
	{$[4; 5)$}
	{\True $[2; 3)$}
	{$[3; 4)$}
	{$[1; 2)$}
	\loigiai{
		tần số của nhóm $[4; 5)$ là lớn nhất bằng $15$.
	}
\end{ex}
\begin{ex}%[2D3V1-4]
	Một nhà khoa học thực hiện một nghiên cứu về độ dài của $30$ con cá trong một hồ cá (đơn vị: centimet). Kết quả thu được được biểu diễn trong bảng sau:
	\begin{center}
		\begin{tabular}{|c|c|c|}
			\hline 
			Nhóm & Tần số & Tần số tích luỹ \\
			\hline 
			$[10; 15)$ & $4$ & $4$ \\
			$[15; 20)$ & $8$ & $12$ \\
			$[20; 25)$ & $6$ & $18$ \\
			$[25; 30)$ & $9$ & $27$ \\
			$[30; 35)$ & $3$ & $30$ \\
			\hline 
			& $n=30$ & \\
			\hline
		\end{tabular}
	\end{center}	
	Mệnh đề nào sau đây là \textbf{sai}?
	\choice
	{Cỡ mẫu $n=30$}
	{\True $Q_1=18~(\mathrm{cm})$}
	{$Q_3=29{,}5~(\mathrm{cm})$}
	{Khoảng tứ phân vị $\Delta Q =15~(\mathrm{cm})$}	
	\loigiai{
		Số lượng mẫu là $n=30$.\\
		Ta có $\dfrac{n}{4}=\dfrac{30}{4}=7{,}5$, và nằm trong nhóm $[15; 20)$. \\
		Vậy $Q_1=15+\left(\dfrac{7{,}5-4}{8}\right) \cdot 5 = 15+\dfrac{3{,}5}{8} = 15{,}44~(\mathrm{cm})$. \\
		Tương tự, ta có $\dfrac{3n}{4}=\dfrac{3 \times 30}{4}=22{,}5$, và nằm giữa nhóm $[25; 30)$. \\
		Vậy $Q_3=25+\left(\dfrac{22{,}5-18}{9}\right) \cdot 5 = 25+\dfrac{4{,}5}{9} = 29{,}5~(\mathrm{cm})$. \\
		Khoảng tứ phân vị của mẫu số liệu là $\Delta Q = Q_3 - Q_1 = 29{,}5 - 15{,}44 = 14{,}06~(\mathrm{cm})$.
	}
\end{ex}
\begin{ex}%[2D3V1-3]
	Thống kê số giờ học thêm mỗi tuần của các học sinh trong một lớp trong $30$ ngày được cho trong bảng sau:
	\begin{center}
		\begin{tabular}{|c|c|c|c|c|c|}
			\hline
			Số giờ học thêm (giờ) & $\left[ 1;3\right)$ & $\left[ 3;5\right)$ & $\left[ 5;7\right)$ & $\left[ 7;9\right)$ & $\left[ 9;11\right)$ \\
			\hline
			Số ngày & $5$ & $8$ & $9$ & $6$ & $2$ \\
			\hline
		\end{tabular}
	\end{center}
	Tứ phân vị thứ ba của mẫu số liệu trên là
	\choice
	{$Q_3=7{,}191$}
	{$Q_3=8{,}167$}
	{\True $Q_3=7{,}167$}
	{$Q_3= 8{,}191$}
	\loigiai{
		Cỡ mẫu: $n=30$.\\
		Gọi $x_1, x_2, \ldots, x_{30}$ là mẫu số liệu gốc gồm số giờ học thêm mỗi tuần của các học sinh trong lớp được sắp xếp theo thứ tự không giảm.\\
		Ta có $x_1, x_2, x_3, x_4, x_5 \in \left[ 1;3\right)$ \\
		$x_6, x_7, x_8, x_9, x_{10}, x_{11}, x_{12}, x_{13} \in \left[ 3;5\right)$ \\
		$x_{14}, x_{15}, x_{16}, x_{17}, x_{18}, x_{19}, x_{20}, x_{21}, x_{22} \in \left[ 5;7\right)$ \\
		$x_{23}, x_{24}, x_{25}, x_{26}, x_{27}, x_{28} \in \left[ 7;9\right)$ \\
		$x_{29}, x_{30} \in \left[ 9;11\right)$ \\
		Tứ phân vị thứ ba của mẫu số liệu gốc là $\dfrac{1}{2}(x_{23} + x_{24}) \in \left[ 7;9\right)$. Do đó, tứ phân vị thứ ba của mẫu số liệu ghép nhóm là
		\begin{align*}
			Q_3 = 7 + \left(\dfrac{3 \cdot \dfrac{30}{4} - 22}{6}\right) \cdot 2 = 7 + \left(\dfrac{22{,}5 - 22}{6}\right) \cdot 2 = 7 + \dfrac{0{,}5}{6} \cdot 2 = 7 + 0{,}167 \approx 7{,}167.
		\end{align*}
		Vậy tứ phân vị thứ ba cần tìm là: $ Q_3 = 7{,}167$.
	}
\end{ex}
\begin{ex}%[2D3H1-3]
	Một nhà tâm lý học thực hiện một nghiên cứu về số lượng tin nhắn mà 15 cặp đôi yêu nhau gửi cho nhau trong một tuần (đơn vị: tin nhắn). Kết quả thu được được biểu diễn trong bảng sau:
	\begin{center}
		\begin{tabular}{|c|c|c|}
			\hline 
			Nhóm & Tần số & Tần số tích luỹ \\
			\hline 
			$[20; 25)$ & $3$ & $3$ \\
			$[25; 30)$ & $4$ & $7$ \\
			$[30; 35)$ & $2$ & $9$ \\
			$[35; 40)$ & $3$ & $12$ \\
			$[40; 45)$ & $3$ & $15$ \\
			\hline 
			& $n=15$ & \\
			\hline
		\end{tabular}
	\end{center}
	Tứ phân vị thứ nhất là
	\choice
	{$Q_1=24{,}9375$}
	{$Q_1=24{,}9382$}
	{\True $Q_1=25{,}9375$}
	{$Q_1= 25{,}9382$}
	\loigiai{
		Số lượng mẫu là $n=15$.\\
		Ta có $\dfrac{n}{4}=\dfrac{15}{4}=3{,}75$, và nằm trong nhóm $[25; 30)$. \\
		Vậy $Q_1=25+\left(\dfrac{3{,}75-3}{4}\right) \cdot 5 = 25+\left(\dfrac{0{,}75}{4}\right) \cdot 5 = 25+0{,}9375 = 25{,}9375$.
	}
\end{ex}
\begin{ex}%[2D3N1-1]
	Khoảng tứ phân vị trong thống kê là gì và được sử dụng trong trường hợp nào?
	\choice
	{\True Một phép đo mô tả sự phân bố của một tập dữ liệu bằng cách chia tập dữ liệu thành bốn phần bằng nhau}
	{Một phép đo mô tả sự biến động của dữ liệu bằng cách chia tập dữ liệu thành bốn phần bằng nhau}
	{Một phép đo mô tả sự phân bố của một tập dữ liệu bằng cách chia tập dữ liệu thành hai phần, mỗi phần chứa một phần tư ($25\%$) của dữ liệu}
	{Một phép đo mô tả sự biến động của dữ liệu bằng cách chia tập dữ liệu thành hai phần, mỗi phần chứa một phần tư ($25\%$) của dữ liệu}
	\loigiai{
	Khoảng tứ phân vị là một phép đo mô tả sự phân bố của một tập dữ liệu bằng cách chia tập dữ liệu thành bốn phần bằng nhau, mỗi phần chứa một phần tư ($25\%$) của dữ liệu. 
	}
\end{ex}
\Closesolutionfile{ans}
\indapan{6}{ans/ans-2-B9-De2-LC}

\cauds
\Opensolutionfile{ans}[ans/ans-2-B9-De2-DS]
\begin{ex}%[2D3H1-4]%Câu 1
	Bạn Trang thống kê lại chiều cao (đơn vị: cm) của các bạn học sinh nữ lớp 12A và lớp 12B ở bảng sau.
	\begin{center}
		\begin{tabular}{|c|c|c|c|c|c|c|}
		\hline Lớp Chiều cao &$[155; 160)$&$[160; 165)$&$[165; 170)$&$[170; 175)$&$[175; 180)$&$[180; 185)$\\
		\hline 12A & $4$ & $9$ & $12$ & $9$ & $0$ & $1$\\
		\hline 12B & $8$ & $12$ & $8$ & $5$ & $2$ & $0$\\
		\hline
	\end{tabular}
	\end{center}
	\choiceTF
	{\True Cỡ mẫu của 12B là $n=35$}
	{\True Tứ phân vị thứ nhất của mẫu số liệu lớp 12A là $\dfrac{5855}{36}$}
	{Tứ phân vị thứ ba của mẫu số liệu lớp 12B là $\dfrac{3845}{24}$}
	{Mẫu số liệu của 12A ít phân tán hơn 12B}
	\loigiai{
	\begin{enumerate}[1.]
		\item Xét ở lớp 12A.\\
		Ta có cỡ mẫu $n=35$.\\
		Gọi $x_1; x_2; \ldots; x_{35}$ là mẫu số liệu gốc gồm chiều cao của $35$ bạn nữ lớp 12A.\\
		Ta có 
		\begin{itemize}
			\item $x_1$,$\ldots$,$x_4 \in[155; 160)$;
			\item $x_5$,$\ldots$,$x_{13} \in[160; 165)$;
			\item $x_{14}$,$\ldots$,$x_{25}\in[165; 170)$;
			\item $x_{26}$, $\ldots$, $x_{34}\in[170; 175)$;
			\item $x_{35}\in$ $[180; 185)$.
		\end{itemize}    
		Tứ phân vị thứ nhất của mẫu số liệu lớp $12A$ là $\dfrac{x_8+x_9}{2}\in[160; 165)$. Do đó, tứ phân vị thứ nhất của mẫu số liệu lớp $12A$ là
		$$Q_1=160+\dfrac{\dfrac{35}{4}-4}{9}\cdot(165-160)=\dfrac{5855}{36}.$$
		Tứ phân vị thứ ba của mẫu số liệu lớp $12A$ là $\dfrac{x_{26}+x_{27}}{2}\in[170; 175)$. Do đó, tứ phân vị thứ ba của mẫu số liệu lớp $12A$ là
		$$Q_3=170+\dfrac{\dfrac{3 \cdot 35}{4}-(4+9+12)}{9}\cdot(175-170)=\dfrac{6145}{36}.$$
		Vậy khoảng tứ phân vị của mẫu số liệu lớp $12A$ là $\Delta_Q=Q_3-Q_1=\dfrac{290}{36} \approx 8{,}06$.
		\item Xét ở lớp $12B$.\\
		Ta có cỡ mẫu $n=35$.\\
		Gọi $x_1 $; $x_2 $;$\ldots $;$x_{35}$ là mẫu số liệu gốc gồm chiều cao của $35$ bạn nữ lớp $12B$.\\
		Ta có
		\begin{itemize}
			\item $x_1$,$\ldots$,$ x_5 \in[155; 160)$;
			\item $x_6$,$\ldots$,$ x_{14} \in[160; 165)$;
			\item $x_{15}$,$\ldots$,$ x_{22} \in[165; 170)$;
			\item $x_{23}$, $x_{24} \in[170; 175)$;
			\item $x_{35} \in[175; 180)$.
		\end{itemize}
		Tứ phân vị thứ nhất của mẫu số liệu lớp $12B$ là $\dfrac{x_8+x_9}{2}\in[160; 165)$. Do đó, tứ phân vị thứ nhất của mẫu số liệu lớp $12B$ là
		$$Q_1=160+\dfrac{\dfrac{35}{4}-8}{12}\cdot(165-160)=\dfrac{3845}{24}.$$
		Tứ phân vị thứ ba của mẫu số liệu lớp $12B$ là $\dfrac{x_{26}+x_{27}}{2}\in[165; 170)$. Do đó, tứ phân vị thứ ba của mẫu số liệu lớp $12B$ là
		$$Q_3=165+\dfrac{\dfrac{3 \cdot 35}{4}-(8+12)}{8}\cdot(170-165)=\dfrac{5305}{32}.$$
		Vậy khoảng tứ phân vị của mẫu số liệu lớp $12B$ là $\Delta_Q=Q_3-Q_1=\dfrac{535}{96} \approx 5{,}57$.\\
		Vậy mẫu số liệu của lớp 12B ít phân tán hơn mẫu số liệu của lớp $12A$.
	\end{enumerate}
	\begin{itemchoice}
		\itemch \textbf{Đúng.}
		\itemch \textbf{Đúng.}
		\itemch \textbf{Sai.}
		\itemch \textbf{Sai.}
	\end{itemchoice}
}
\end{ex}

\begin{ex}%[2D3H1-4]%Câu 2
	Hằng ngày ông Thắng đều đi xe buýt từ nhà đến cơ quan. Dưới đây là bảng thống kê thời gian của 100 lần ông Thắng đi xe buýt từ nhà đến cơ quan. Xét tính đúng - sai của các mệnh đề sau.
	\begin{center}
		\begin{tabular}{|c|c|c|c|c|c|c|}
		\hline Thời gian (phút) &$[15; 18)$&$[18; 21)$&$[21; 24)$&$[24; 27)$&$[27; 30)$&$[30; 33)$\\
		\hline Số lượt & $22$ & $38$ & $27$ & $8$ & $4$ & $1$\\
		\hline
	\end{tabular}
	\end{center}
	\choiceTF
	{Cỡ mẫu $n=33$}
	{\True $Q_1=\dfrac{693}{38}$}
	{\True $\Delta_Q=\dfrac{505}{114}$}
	{\True Đi hết $31$ phút là một giá trị ngoại lệ của mẫu số liệu ghép nhóm}
	\loigiai{
	Cỡ mẫu $n=100$.\\
	Gọi $x_1 $; $x_2 $;$\ldots $;$x_{100}$ là mẫu số liệu gốc gồm thời gian $100$ lần đi xe buýt của ông Thắng.\\
	Ta có 
	\begin{itemize}
		\item $x_1$,$\ldots$,$ x_{22}\in[15 ; 18)$;
		\item $x_{23}$,$\ldots$,$ x_{60}\in[18 ; 21)$;
		\item $x_{61}$,$\ldots$,$ x_{87}\in[21 ; 24)$;
		\item $x_{88}$,$\ldots$,$ x_{95}\in[24 ; 27)$;
		\item $x_{96}$,$\ldots$,$ x_{99}\in[27 ; 30)$;
		\item $x_{100}\in[30 ; 33)$.
	\end{itemize}
	Tứ phân vị thứ nhất của mẫu số liệu gốc là $\dfrac{x_{25}+x_{26}}{2}\in[18 ; 21)$. Do đó, tứ phân vị thứ nhất của mẫu số liệu ghép nhóm là
	$$Q_1=18+\dfrac{\dfrac{100}{4}-22}{38}\cdot(21-18)=\dfrac{693}{38}.$$
	Tứ phân vị thứ ba của mẫu số liệu gốc là $\dfrac{x_{75}+x_{76}}{2}\in[21 ; 24)$. Do đó, tứ phân vị thứ ba của mẫu số liệu ghép nhóm là
	$$Q_3=21+\dfrac{\dfrac{3 \cdot 100}{4}-(22+38)}{27}\cdot(24-21)=\dfrac{68}{3}.$$
	Khoảng tứ phân vị của mẫu số liệu ghép nhóm là $\Delta_Q=Q_3-Q_1=\dfrac{505}{114}$.\\
	Khi ông Thắng đi hết $31$ phút, thời gian đi của ông thuộc nhóm $[30 ; 33)$.\\
	Vì $Q_3+1,5 \Delta_Q=\dfrac{6683}{228}<30$ nên thời gian của lần ông Thắng đi hết $31$ phút là giá trị ngoại lệ của mẫu số liệu ghép nhóm.
	\begin{itemchoice}
		\itemch \textbf{Sai.}
		\itemch \textbf{Đúng.}
		\itemch \textbf{Đúng.}
		\itemch \textbf{Đúng.}
	\end{itemchoice}
}
\end{ex}

\begin{ex}%[2D3H1-4]%Câu 3	
	Cho bảng biểu diễn mẫu số liệu ghép nhóm về chiều cao của $50$ mẫu cây ở một vườn thực vật (đơn vị centimet). Xét tính đúng sai của các mệnh đề sau.
	\begin{center}
		\begin{tabular}{|c|c|c|}
		\hline Nhóm & Tần số & Tần số tích luỹ \\
		\hline $[40; 45)$ & $7$ & $7$ \\
		$[45; 50)$& $12$ & $19$ \\
		$[50; 55)$& $9$ & $28$ \\
		$[55; 60)$& $10$ & $38$ \\
		$[60; 65)$& $8$ & $46$ \\
		$[65; 70)$& $4$ & $50$ \\
		\hline & $n=50$ & \\
		\hline
	\end{tabular}
	\end{center}
	\choiceTF
	{\True Cỡ mẫu $n=50$}
	{$Q_1=45$}
	{$Q_3=67{,}5$}
	{$\Delta_Q\approx 54{,}075$}
	\loigiai{
	Số phần tử của mẫu là $n=50$.\\
	Ta có $\dfrac{n}{4}=\dfrac{42}{4}=10{,}5$ mà $5<10{,}5<15$.
	$$Q_1=45+\left(\dfrac{10{,}5-5}{10}\right) \cdot 5=47{,}75~(\mathrm{cm}).$$
	Ta có $\dfrac{3 n}{4}=\dfrac{3 \cdot 42}{4}=31{,}5$ mà $31<31{,}5<38$.\\
	Áp dụng công thức, ta có tứ phân vị thứ ba là
	$$Q_3=60+\left(\dfrac{31{,}5-31}{7}\right) \cdot 5 \approx 60{,}4~(\mathrm{cm}).$$
	Vậy khoảng tứ phân vị của mẫu số liệu ghép nhóm đã cho là
	\[
	\Delta_Q=Q_3-Q_1 \approx 60{,}4-47{,}75 \approx 12{,}65~(\mathrm{cm}).
	\]
		\begin{itemchoice}
			\itemch \textbf{đúng.}
			\itemch \textbf{Sai.}
			\itemch \textbf{Sai.}
			\itemch \textbf{Sai.}
		\end{itemchoice}}
\end{ex}

\begin{ex}%[2D3H1-4]%Câu 4
	Ở một phòng điều trị nội trú của bệnh viện, dữ liệu thống kê thời gian ngủ hằng đêm của hai bệnh nhân trong suốt một tháng được tổng hợp bởi hai bảng dưới đây:
	\begin{center}
		\begin{tabular}{cc}
		Bảng a. Thời gian ngủ của bệnh nhân $A$
		&
		Bảng b. Thời gian ngủ của bệnh nhân $B$\\
		\begin{tabular}{|c|c|}
			\hline Thời gian (phút) & Số đêm (tần số) \\
			\hline $[180; 240)$ & $5$ \\
			\hline $[240; 300)$ & $5$ \\
			\hline $[300; 360)$ & $10$ \\
			\hline $[360; 420)$ & $6$ \\
			\hline $[420; 480)$ & $4$ \\
			\hline
		\end{tabular}  
		&
		\begin{tabular}{|c|c|}
			\hline Thời gian (phút) & Số đêm (tần số) \\
			\hline $[180; 240)$ & $2$ \\
			\hline $[240; 300)$ & $9$ \\
			\hline $[300; 360)$ & $12$ \\
			\hline $[360; 420)$ & $5$ \\
			\hline $[420; 480)$ & $2$ \\
			\hline
		\end{tabular}  
	\end{tabular}
	\end{center}
	\choiceTF
	{Kích thước mẫu số liệu là $45$}
	{\True $Q_1A=270$}
	{$\Delta_QB=81$}
	{Bệnh nhân $A$ ngủ ổn định hơn bệnh nhân $B$}
	\loigiai{
	Kích thước của hai mẫu số liệu đều là $n=30$. Ta có $\dfrac{N}{4}=\dfrac{15}{2}=7{,}5$; $\dfrac{3 N}{4}=\dfrac{45}{2}=22{,}5$.
	Đối với mẫu số liệu về thời gian ngủ của bệnh nhân $A$, ta tìm các tứ phân vị $Q_1A$, $Q_3A$ và khoảng tứ phân vị $\Delta_QA$ qua bảng tần số tích luỹ dưới đây:
	\begin{center}
		\begin{tabular}{|c|c|c|}
		\hline Thời gian (phút) & Tần số & \begin{tabular}{c}
			Tần số \\
			tích luỹ
		\end{tabular}\\
		\hline $[180; 240)$ & $5$ & $5$ \\
		\hline $[240; 300)$ & $5$ & $10$ \\
		\hline $[300; 360)$ & $10$ & $20$ \\
		\hline $[360; 420)$ & $6$ & $26$ \\
		\hline $[420; 480)$ & $4$ & $30$ \\
		\hline
	\end{tabular}
	\end{center}
	Nhóm chứa $Q_1A$ là $[240; 300)$ nên $Q_1A=240+\dfrac{7,5-5}{5}\cdot 60=270$.\\
	Nhóm chứa $Q_3A$ là $[360; 420)$ nên $Q_3A=360+\dfrac{22,5-20}{6}\cdot 60=385$.\\
	Suy ra $\Delta_QA=385-270=115$.
	Đối với mẫu số liệu về thời gian ngủ của bệnh nhân $B$, ta tìm các tứ phân vị $Q_1B$, $Q_3B$ và khoảng tứ phân vị $\Delta_QB$ qua bảng tần số tích luỹ dưới đây:
	\begin{center}
			\begin{tabular}{|c|c|c|}
			\hline Thời gian (phút) & Tần số & \begin{tabular}{c}
				Tần số \\
				tích luỹ
			\end{tabular}\\
			\hline $[180; 240)$ & $2$ & $2$ \\
			\hline $[240; 300)$ & $9$ & $11$ \\
			\hline $[300; 360)$ & $12$ & $23$ \\
			\hline $[360; 420)$ & $5$ & $28$ \\
			\hline $[420; 480)$ & $2$ & $30$ \\
			\hline
		\end{tabular}
	\end{center}
	Nhóm chứa $Q_1B$ là $[240; 300)$ nên $Q_1B=240+\dfrac{7,5-2}{9}\cdot 60 \approx 276{,}67$.\\
	Nhóm chứa $Q_3B$ là $[300; 360)$ nên $Q_3B=300+\dfrac{22,5-11}{12}\cdot 60=357{,}5$.\\
	Suy ra $\Delta_QB \approx 357,5-276,67 \approx 80{,}83$.\\
	Vì $\Delta_QA>\Delta_QB$ nên bệnh nhân $B$ có thời gian ngủ ổn định hơn.
	\begin{itemchoice}
		\itemch \textbf{Đúng.}
		\itemch \textbf{Sai.}
		\itemch \textbf{Sai.}
		\itemch \textbf{Sai.}
	\end{itemchoice}
}
\end{ex}
\Closesolutionfile{ans}
\indapan{3}{ans/ans-2-B9-De2-DS}

\Opensolutionfile{ans}[ans/ans-2-B9-De2-KQ]
\caukq
\begin{ex}%[2D3H1-2]
	Bảng dưới đây biểu diễn mẫu số liệu ghép nhóm về số lần khách hàng ghé thăm một quán cà phê trong một tháng.
	\begin{center}
		\begin{tabular}{|c|c|c|c|c|c|}
			\hline Số lần & $[5 ; 10)$ & $[10 ; 15)$ & $[15 ; 20)$ & $[20 ; 25)$ & $[25 ; 30)$ \\
			\hline Số khách & $4$ & $8$ & $12$ & $20$ & $16$ \\
			\hline
		\end{tabular}
	\end{center}
	Khoảng biến thiên của mẫu số liệu ghép nhóm trên bằng bao nhiêu?
	\shortans[]{$25$}
	\loigiai{
		Trong mẫu số liệu ghép nhóm ở bảng trên, ta có đầu mút trái của nhóm đầu tiên là $a_1=5$, và đầu mút phải của nhóm cuối cùng là $a_6=30$.\\
		Vậy khoảng biến thiên của mẫu số liệu ghép nhóm đó là\\
		\[R = a_6 - a_1 = 30 - 5 = 25\].
	}
\end{ex}
\begin{ex}%[2D3V1-3]
	Bảng sau biểu diễn mẫu số liệu ghép nhóm về số phút mà $50$ người sử dụng dịch vụ internet trong một ngày.
	\begin{center}
		\begin{tabular}{|c|c|c|c|c|c|}
			\hline Nhóm  & $[20 ; 30)$ & $[30 ; 40)$ & $[40 ; 50)$ & $[50 ; 60)$ & $[60 ; 70)$ \\
			\hline Tần số & $5$ & $12$ & $15$ & $10$ & $8$ \\
			\hline
		\end{tabular}
	\end{center}
	Khoảng tứ phân vị của mẫu số liệu ghép nhóm trên bằng bao nhiêu (làm tròn đến $1$ chữ số thập phân)?
	\shortans[]{$19{,}2$}
	\loigiai{
		Ta có số phần tử của mẫu là $n=50$.\\
		Áp dụng công thức, ta có tứ phân vị thứ nhất là 
		\[
		Q_1 = 30 + \left( \dfrac{12{,}5 - 5}{12} \right) \cdot 10 = 30 + \left( \dfrac{7{,}5}{12} \right) \cdot 10 = 30 + 6{,}25 = 36{,}3~\text{(phút)}.
		\]
		Tứ phân vị thứ ba là 
		\[
		Q_3 = 50 + \left( \dfrac{37{,}5 - 32}{10} \right) \cdot 10 = 50 + \left( \dfrac{5{,}5}{10} \right) \cdot 10 = 50 + 5{,}5 = 55{,}5~\text{(phút)}.
		\]
		Vậy khoảng tứ phân vị của mẫu số liệu ghép nhóm đã cho là 
		\[
		\Delta_Q = Q_3 - Q_1 = 55{,}5 - 36{,}3 = 19{,}2~\text{(phút)}.
		\]
	}
\end{ex}
\begin{ex}%[2D3V1-3]
	Bảng sau biểu diễn mẫu số liệu ghép nhóm về tốc độ (đơn vị: km/h) của các phương tiện giao thông trên đường cao tốc.
	\begin{center}
		\begin{tabular}{|c|c|c|c|c|c|}
			\hline Nhóm  & $[60 ; 70)$ & $[70 ; 80)$ & $[80 ; 90)$ & $[90 ; 100)$ & $[100 ; 110)$ \\
			\hline Tần số & $10$ & $20$ & $30$ & $25$ & $15$ \\
			\hline
		\end{tabular}
	\end{center}
	Tính khoảng tứ phân vị của mẫu số liệu ghép nhóm trên (làm tròn đến 1 chữ số thập phân).
	\shortans[]{$18{,}5$}
	\loigiai{
		Ta có số phần tử của mẫu là $n=100$.\\
		Áp dụng công thức, ta có tứ phân vị thứ nhất là 
		\[
		Q_1 = 70 + \left( \dfrac{25 - 10}{20} \right) \cdot 10 = 70 + \left( \dfrac{15}{20} \right) \cdot 10 = 70 + 7{,}5 = 77{,}5~\text{(km/h)}.
		\]
		Tứ phân vị thứ ba là 
		\[
		Q_3 = 90 + \left( \dfrac{75 - 60}{25} \right) \cdot 10 = 90 + \left( \dfrac{15}{25} \right) \cdot 10 = 90 + 6 = 96~\text{(km/h)}.
		\]
		Vậy khoảng tứ phân vị của mẫu số liệu ghép nhóm đã cho là 
		\[
		\Delta_Q = Q_3 - Q_1 = 96 - 77{,}5= 18{,}5~\text{(km/h)}.
		\]
	}
\end{ex}
%%%%%%%%%%%%%%%
\begin{ex}%[2D3N1-2]
	Cho bảng biểu diễn mẫu số liệu ghép nhóm về độ tuổi của cư dân trong một khu phố.
	\begin{center}    
		\begin{tabular}{|c|c|c|c|c|c|c|}
			\hline 
			Nhóm & $[20 ; 30)$& $[30 ; 40)$ & $[40 ; 50)$ & $[50 ; 60)$& $[60 ; 70)$ & $[70 ; 80)$ \\ 
			\hline 
			Tần số & $25$ &$20$&$20$ &$15$&$14$&$6$ \\
			\hline
		\end{tabular}
	\end{center}
	Tính khoảng biến thiên của mẫu số liệu ghép nhóm đó.
	\shortans{$60$}
	\loigiai{
		Trong mẫu số liệu ghép nhóm ở bảng $10$, ta có đầu mút trái của nhóm $1$ là $a_1=20$, đầu mút phải của nhóm $6$ là $a_7=80$.\\
		Vậy khoảng biến thiên của mẫu số liệu ghép nhóm đó là $R=a_7-a_1=80-20=60$.
	}
\end{ex}
\begin{ex}%[2D3V1-3]
	Cho bảng biểu diễn mẫu số liệu thống kê thu nhập của các cặp vợ chồng trong một khu phố (đơn vị: triệu đồng).
	\begin{center}
		\begin{tabular}{|c|c|c|c|c|c|c|}
			\hline
			Nhóm & $[10 ; 15)$ & $[15 ; 20)$ & $[20 ; 25)$ & $[25 ; 30)$ & $[30 ; 35)$ & $[35 ; 40)$ \\
			\hline
			Tần số & $12$ & $20$ & $15$ &$8$ &$4$ &$1$ \\
			\hline
		\end{tabular}
	\end{center}
	Tính khoảng tứ phân vị của mẫu số liệu thu nhập của các cặp vợ chồng đã cho(làm tròn đến hàng phần trăm).
	\shortans[]{$8{,}58$}
	\loigiai{
		Số lượng cặp vợ chồng trong mẫu là $n = 60$.\\
		Ta có $\frac{n}{4} = \frac{60}{4} = 15$, mà $15 < 32$.\\
		Áp dụng công thức, ta có tứ phân vị thứ nhất là
		$Q_1 = 15 + \left(\frac{15 - 12}{20}\right) \cdot 5 = 15{,}75 ~\text{(triệu đồng)}$.\\
		Ta có $\frac{3n}{4} = \frac{3 \cdot 60}{4} = 45$, mà $32 < 45 < 47$.\\
		Áp dụng công thức, ta có tứ phân vị thứ ba là
		$Q_3 = 20 + \left(\frac{45 - 32}{15}\right) \cdot 5\approx 24{,}33 ~\text{(triệu đồng)}$.\\
		Vậy khoảng tứ phân vị của mẫu số liệu thu nhập của các cặp vợ chồng đã cho là\\
		$\Delta_Q = Q_3 - Q_1 = 24{,}33 - 15{,}75 = 8{,}58 ~\text{(triệu đồng)}$.
	}
\end{ex}
%
\begin{ex}%[2D3V1-3]
	Thu nhập của các cặp vợ chồng trong một năm ở một khu phố được thể hiện qua biểu đồ cột sau.
	\begin{center}
		\begin{tikzpicture}[scale=1, font=\footnotesize, line join=round, line cap=round, >=stealth]
			\coordinate [label= left: $0$](O) at (0,0) ;
			\coordinate [label= left: Số cặp vợ chồng](y) at (0,8) ;
			\coordinate [label= below left: Thu nhập (triệu đồng)](x) at (10,0) ;
			\foreach \i in {2,3,4,5,6,7}
			{\draw[-,color=gray!40] (0,\i)--(10,\i); }
			\draw (0,1)node[left]{$10$}(0,2)node[left]{$20$}(0,3)node[left]{$30$}(0,4)node[left]{$40$}(0,5)node[left]{$50$}(0,6)node[left]{$60$}(0,7)node[left]{$70$};
			%--------------------------------------------
			\foreach \i/\j/\a in 
			{0.25/0.2/4,0.75/1.0/20,1.25/2.5/50,1.75/3.0/60,2.25/2.0/40,2.75/1.5/30,3.25/1.0/20,
				3.75/0.5/10,4.25/0.4/8,4.75/0.2/4,5.25/0.1/2}{
				\draw[] ($(\i,\j)+(0,0.2)$) node[scale=0.7]{$\a$};
			}
			%-----------------------------------------------
			\foreach \i/\j in 
			{0/0.2,0.5/1.0,1/2.5,1.5/3.0,2/2.0,2.5/1.5,3/1.0,
				3.5/0.5,4/0.4,4.5/0.2,5/0.1}{
				\fill[blue!20] (\i,0) rectangle+ (0.5,\j);
				\draw[] (\i,0) rectangle+ (0.5,\j);
			}
			\foreach \i/\j in 
			{0.25/[401-450],0.75/[451-500],1.25/[501-550],1.75/[551-600],2.25/[601-650],2.75/[651-700],3.25/[701-750],
				3.75/[751-800],4.25/[801-850],4.75/[851-900],5.25/[901-950]}{
				\draw($(\i,0)+(0,-1)$) node[rotate=90]{\j};
			}   
			\foreach \i/\j in {}{\draw (\i,\j) node[above]{\j};}           
			\draw[->] (O)--(x);
			\draw[->] (O)--(y);
		\end{tikzpicture}
	\end{center}
	Tính Khoảng tứ phân vị của mẫu thu nhập trên (làm tròn đến hàng đơn vị).
	\shortans[]{$94$}
	\loigiai{
		Ta thống kê theo bảng tần số như sau
		\begin{center}
			\begin{tabular}{|c|c|c|c|c|c|c|}
				\hline Nhóm  &$[401 ; 450]$ & $[451 ; 500]$ & $[501 ; 550]$ &	$[551 ; 600]$ &	$[601 ; 650]$ & $[651 ; 700]$ \\
				\hline Tần số & $4$ & $20$ & $50$ & $60$ &$40$& $30$  \\
				\hline
			\end{tabular}
		\end{center}
		\begin{center}
			\begin{tabular}{|c|c|c|c|c|c|}
				\hline Nhóm  & $[701 ; 750]$ &$[751 ; 800]$&$[801 ; 850]$& $[851 ; 900]$ & $[901 ; 950]$ \\
				\hline Tần số & $20$ &$10$& $8$ & $4$ & $2$ \\
				\hline
			\end{tabular}
		\end{center}
		Cỡ mẫu của bảng số liệu là $n=200$ cặp vợ chồng.\\
		Tứ phân vị thứ nhất là $Q_1=501+\dfrac{\dfrac{200}{4}-24}{50}\cdot 50=527$.\\
		Tứ phân vị thứ ba là $Q_3=601+\dfrac{\dfrac{3}{4}\cdot 200-134}{40}\cdot 50=621$.\\
		Do đó, khoảng tứ phân vị là $\Delta Q=621 - 560 = 527 = 94$.
		
	}
\end{ex}
\Closesolutionfile{ans}
\indapan{6}{ans/ans-2-B9-De2-KQ}